\documentclass[10pt]{article}
\usepackage[margin=20mm]{geometry}
\usepackage{amsmath,amssymb,mathtools,fontspec,unicode-math,xurl,hyperref}
\IfFontExistsTF{DejaVu Serif}{\setmainfont{DejaVu Serif}}{\setmainfont{Latin Modern Roman}}
\setmathfont{latinmodern-math.otf}
\setlength{\parindent}{0pt}\setlength{\parskip}{6pt}\sloppy
\title{Technical review of Revisiting mixed geometry}
\author{AI-assisted technical review}
\date{}
\begin{document}
\maketitle
Abstract: The exact final-author arXiv v4 manuscript is reviewed through its relative-tensor construction, enriched Hom formula, weight and perverse structures, and Hecke comparison. The unrestricted constructible pushforward statement fails for \(B\mathbb G_m\to\mathrm{pt}\) and is assessed as Major in that scope. Four local errors concern continuity of adjoints, Galois conjugates, the covering condition for descent, and the eigenvalue index set of an optional variant. The central renormalized construction and the Hecke equivalence are not disproved. Actual source and report image inspection remains blocked.\par
\section*{Exact source and scope}
The authors are Quoc P. Ho and Penghui Li. The source is the 63-page arXiv:2202.04833v4, associated with DOI 10.4171/jems/1739 (https://doi.org/10.4171/jems/1739) [paper, publisher]. Its arXiv stamp is 1 December 2025; its internal first-page date is 2 December 2025. The arXiv record labels v4 final. A different date in a subsequently retrieved HTML search result, and an author’s separately hosted PDF, are not used to establish equality with the bound file. All references below are to its actual PDF/article pages.\par
All 63 extracted text pages were actually read. The checks reconstruct selected implications rather than independently rebuilding the entire higher-categorical and mixed-sheaf foundations. Throughout, the actual convention on page 4 is finite-type Artin stacks with affine stabilizers; the introductory phrase arbitrary Artin stacks is not read as removing that convention.\par
The three categories \(\mathsf{Shv}(Y)\), \(\mathsf{Shv}(Y)^{\mathsf{ren}}\), and \(\mathsf{Shv}_{\mathsf{gr},\mathsf c}(Y)\) have different roles. Constructible objects are compact by definition in the renormalized category, not necessarily in the usual category. On a point, \(\mathsf{Shv}_{\mathsf{gr}}(\mathrm{pt})=\mathsf{Vect}^{\mathsf{gr}}\) and \(\mathsf{Shv}_{\mathsf{gr},\mathsf c}(\mathrm{pt})=(\mathsf{Vect}^{\mathsf{gr}})^{\mathsf c}\). A compact graded complex has finitely many nonzero grading components, each perfect. This distinction is essential to the first finding.\par
\section*{P1: constructible pushforward is overstated}
\subsection*{A counterexample within the stated stack class}
Proposition 3.5.5(ii), page 40, is stated for a morphism \(f:Y\to Z\) in \(\mathsf{Stk}_k\) and writes the pushforwards as functors \[f_* \text{ (respectively }f_!\text{)}:
\mathsf{Shv}_{\mathsf{gr},\mathsf c}(Y)\longrightarrow\mathsf{Shv}_{\mathsf{gr},\mathsf c}(Z).\] The unrestricted \(f_*\) part is false. Let \[f:B\mathbb G_m\longrightarrow\mathrm{pt},\qquad
F=\Lambda_{B\mathbb G_m},\] with \(\Lambda=\overline{\mathbb Q}_\ell\). This stack is finite type and has affine stabilizers. The constant sheaf is constructible, and its graded image belongs to the source small category.\par
Its renormalized pushforward is \[f_{*,\mathsf{ren}}F=C^*_{\mathsf{gr}}(B\mathbb G_m)=\Lambda[\beta],
\qquad \deg_{\mathrm{coh}}\beta=\deg_{\mathsf{gr}}\beta=2.\] This is the same cohomology calculation used on pages 9 and 48, with the Frobenius weight retained as grading. It follows also from the stable projective-space approximations to \(B\mathbb G_m\). There is a nonzero component \[\Lambda[-2j]\langle-2j\rangle\] in every grade \(2j\), for \(j\geq0\). Thus the pushforward is not compact in \(\mathsf{Vect}^{\mathsf{gr}}\) and is not in \(\mathsf{Shv}_{\mathsf{gr},\mathsf c}(\mathrm{pt})\).\par
This failure does not rely only on a characterization quoted from the paper. Let \(R_{\leq N}\) be the finite graded subcomplex spanned by \(1,\beta,\ldots,\beta^N\). Then \(R=\operatorname*{colim}_N R_{\leq N}\). If \(R\) were compact, its identity would factor through some \(R_{\leq N}\). In grade \(2N+2\), every such factorization is zero, whereas the identity is nonzero. This proves noncompactness directly. The reviewer script verifies the finite truncation and grading algebra; the infinite conclusion follows from this exact argument, not a finite numerical extrapolation.\par
\subsection*{Actual impact and repair}
This is assessed as Major for the stated all-morphism constructible-functor conclusion. On this allowed input, the displayed functor into the small category does not exist, so its asserted bounded weight-structure property cannot be evaluated there. Replacing the target by a large category is a substantive change of scope.\par
The earlier formulation in Remark 1.5.2 and Theorem 2.7.12 already recognizes this issue and uses renormalized large categories. Those constructions are not refuted. Nor is the Hecke application: its convolution pushforwards are representable and proper, so the relevant constructibility requirement is satisfied.\par
The small-category statement should be restricted to maps or objects for which pushforward preserves constructibility, using for instance the \(\mathrm{UC}_*\) class defined in Section 2.8. A statement in a large category would instead need a specified extended weight structure and an argument for its claimed exactness. Such a general extension is not supplied by simply applying the bounded small-category result. The counterexample and the necessary restriction are verified here; no general unbounded replacement theorem is claimed.\par
\section*{P2: a right adjoint need not be a strict module functor}
Corollary A.4.7 on page 57 equates laxness and strictness of both members of an adjoint pair over a rigid monoidal category, without requiring continuity of the right adjoint. The immediately preceding Lemma A.4.6 explicitly requires a continuous lax functor.\par
Take \[A=M=N=\mathsf{Vect},\qquad
U=\bigoplus_{j\geq0}\Lambda,\qquad
F=U\otimes-,\qquad G=\operatorname{RHom}(U,-).\] Here \(\mathsf{Vect}\) is compactly generated and rigid in the paper’s sense: its compact objects, not all its objects, are dualizable. The pair \(F\dashv G\) is a genuine adjoint pair between the underlying categories, and \(F\) is a continuous strict \(\mathsf{Vect}\)-module functor. All the printed hypotheses hold.\par
The right-lax comparison obtained by adjunction, evaluated at the scalar object \(U\) and target \(\Lambda\), is \[U\otimes G(\Lambda)=U\otimes U^*
\longrightarrow G(U)=\operatorname{End}(U).\] Its image consists of finite-rank linear maps. The identity of the infinite-dimensional space \(U\) is not in its image. Equivalently, the diagonal infinite matrix belongs to \(\prod_j\bigoplus_i\Lambda\) but not to \(\bigoplus_i\prod_j\Lambda\). Thus \(G\) is not strict.\par
The assessment is Minor in the checked proof chain, because the enriched-Hom proof in A.5 explicitly establishes continuity of its right adjoints using compact inputs, and the renormalized pullback functors preserve compactness. Those essential uses retain the missing hypothesis. The unrestricted corollary, and its unqualified invocation for usual potentially noncontinuous pushforwards in Proposition 2.3.5, still require correction. This assessment is based on the actual uses, not on the location in an appendix.\par
Require \(G\) to be continuous, equivalently require \(F\) to preserve compactness when its source is compactly generated, and then apply Lemma A.4.6. Without this assumption one has the adjoint lax structure; comparison against a dualizable scalar is an equivalence, but comparison against every ind-object need not be. The finite matrix checks retained in the calculation log illustrate the support distinction; the infinite-rank identity above is the proof.\par
\section*{P3: base change produces conjugates}
Lemma 2.7.5 on page 23 writes \(g^*g_*F\simeq F^{\oplus r}\) for finite extension of the base field. The proof of Theorem 3.2.19 on page 35 repeats an identical-copy formula. In general the summands are Galois conjugates: \[g^*g_*F\simeq
\bigoplus_{\sigma\in\operatorname{Gal}(k_m/k_n)}\sigma^*F.\]\par
An explicit example is \(k_n=\mathbb F_2\), \(k_m=\mathbb F_4\), \(Y_n=\mathbb A^1_{\mathbb F_2}\), and \(F=\delta_\alpha\) on \(Y_m\), where \(\alpha^2+\alpha+1=0\). This is a constructible pure perverse skyscraper of weight zero. Its conjugate is supported at \(\alpha^2=\alpha+1\ne\alpha\). Consequently \[g^*g_*\delta_\alpha\simeq\delta_\alpha\oplus\delta_{\alpha^2}
\not\simeq\delta_\alpha^{\oplus2}.\] The stalk at \(\alpha^2\) already distinguishes the two sides. The finite-field arithmetic and support multiplicities were independently computed exactly.\par
This is Minor for the actual arguments. The identity Galois component provides the retract needed in Lemma 2.7.5; full faithfulness has already been established before the idempotent is split. In Theorem 3.2.19 every conjugate is again pure and geometrically simple and therefore belongs to the stated collection of generators. Replacing identical copies by the finite conjugate sum preserves both deductions. No failure of base-extension invariance or generation is inferred from the incorrect stronger formula.\par
\section*{P4 and P5: two auxiliary scope errors}
\subsection*{P4: smoothness alone does not give descent}
Proposition 2.5.10 on page 19 states its equivalence for a smooth morphism \(Z_n\to Y_n\), without a covering condition. Set \[Y_n=\operatorname{Spec}k_n\sqcup\operatorname{Spec}k_n,\qquad
Z_n=\operatorname{Spec}k_n\] and use the inclusion of the first component. This is a smooth map between allowed finite-type schemes. It is a monomorphism, so its entire Čech nerve is constant at \(Z_n\). The asserted functor is the projection \[\mathsf{Vect}^{\mathsf{gr}}\times\mathsf{Vect}^{\mathsf{gr}}\longrightarrow\mathsf{Vect}^{\mathsf{gr}},\] which kills the nonzero sheaf supported on the second point. It cannot be an equivalence, on either the large or the constructible version.\par
This is Minor: require a smooth surjection, or a jointly surjective smooth covering family. That is the hypothesis of the mixed-sheaf descent input used in the proof and of the actual atlas applications. This finding does not assert that smooth-cover descent fails. The separate statement that renormalized sheaves do not satisfy smooth descent is already explicitly made in Remark 2.5.12 and is not a new objection.\par
\subsection*{P5: the optional eigenvalue index excludes an invertible object}
Remark 2.4.6, page 16, describes the index set for \(\mathsf{Vect}^\omega\) as algebraic integers of the permitted Frobenius absolute values. Take \(q=2\), \(n=1\), \(\ell=3\), and the Tate sheaf \(\Lambda(1)\) over \(\mathbb F_2\). It is a continuous rank-one Frobenius representation, pure of weight \(-2\), with eigenvalue \(1/2\). This is an \(\ell\)-adic unit, but it is not an algebraic integer. Its primitive integer minimal polynomial is \(2X-1\); more elementarily, a rational algebraic integer must be an integer.\par
Thus the proposed target omits an allowed tensor-invertible object and is not an appropriate eigenvalue index for the asserted functor. Use the multiplicative group of permitted nonzero Frobenius eigenvalues, with their Weil-weight and continuity conditions, instead of only algebraic integers.\par
This is Minor for the main theory: the optional \(\omega\) variant is not used in its proofs. The principal integer-weight grading explicitly includes negative weights and satisfies \(\operatorname{gr}(F(k))=\operatorname{gr}(F)\langle2k\rangle\). No obstruction to that grading is shown. The full optional variant is not independently developed here.\par
\section*{Checks of the central construction and applications}
\subsection*{Relative tensors and enriched Hom}
The operative definition is \[\mathsf{Shv}_{\mathsf{gr}}(Y)^{\mathsf{ren}}
=\mathsf{Shv}_m(Y_n)^{\mathsf{ren}}\otimes_{\mathsf{Shv}_m(\mathrm{pt}_n)}\mathsf{Vect}^{\mathsf{gr}}.\] The scalar functor records naive Frobenius weights and forgets the Frobenius module structure. It does not take Frobenius invariants. The enriched-Hom formula uses a compact first argument and follows by composing continuous right adjoints in A.5. Those compactness conditions cannot be dropped just because a category is rigid.\par
A point calculation illustrates the distinction. For the pure weight-zero Frobenius matrix \(\left(\begin{smallmatrix}1&1\\0&1\end{smallmatrix}\right)\), the underlying endomorphism space has dimension four, while the commuting endomorphisms have dimension two. The graded Hom formula takes the weight-zero part of the enriched geometric Hom, giving dimension four in this case, rather than its Frobenius-invariant subspace. Exact commutant calculation confirms these dimensions.\par
The finite-field structure map must also be retained: viewing \(\operatorname{Spec}\mathbb F_{q^2}\) over \(\mathbb F_q\) gives a geometrically two-point scheme, whereas viewing it over itself gives a point. Remark 2.4.5 correctly distinguishes these constructions. Invariance under changing a form simultaneously changes the scalar category; it is not the same-base pullback operation of Section 2.5.\par
\subsection*{Weights, perverse objects, and generation}
For pure graded perverse objects of weights \(k,l\), the mixed Hom estimate gives cohomology weights at least \(i+l-k\) in degree \(i\). Taking naive weight zero forces \(i\leq k-l\); the perverse-heart condition gives \(i\geq0\). Distinct weights have no degree-zero maps, while equal weights have Hom concentrated in degree zero. These are the precise orthogonality statements used in the generation and transversality arguments.\par
The proof uses finite weight intervals, the semiorthogonal decomposition lemma, and idempotent closure to generate the small category before invoking Bondarko’s criterion. The corrected Galois-conjugate sum in P3 still supplies the required geometrically simple generators. Mixed sheaf weight filtrations and the external Hom bounds remain inputs.\par
The grading and cohomological shifts have consistent signs: naive grade \(i\) in cohomological degree \(i\) has complex weight zero; \([1]\) raises complex weight by one and \(\langle1\rangle\) lowers it by one. A weight heart need not be a classical category, because negative Hom degrees can survive. The equivalence with bounded chain complexes is therefore correctly conditional on the Hom-purity criterion of Proposition 3.7.2.\par
\subsection*{The target of the graded lift and the Hecke category}
The comparison target is \(\mathsf{Shv}_{\infty,\mathsf c}(Y)\), generated by finite-field mixed descents, not automatically every geometric constructible sheaf. The abelian-heart Ext comparison in Proposition 3.6.9 uses the exact grade-forgetting/constant-grading adjunction on ind-hearts, followed by almost compactness. It is not justified merely by identifying all Ext groups of an arbitrary heart with ambient derived Hom. The stated external almost-compactness result is retained as a dependency.\par
For \(B\backslash G/B\), convolution pulls back along smooth representable maps and pushes forward along proper representable maps. These operations preserve the relevant small categories, so P1’s counterexample does not apply. Hom purity from the cited Bezrukavnikov–Yun results makes the weight heart classical. The polynomial cohomology ring, shearing functors, Bott–Samelson generators, direct summands, and the cited fully-faithfulness calculation then supply the route to the Soergel bimodule equivalence. The shearing functors are only monoidal, as the manuscript explicitly warns; no symmetric-monoidal equivalence is substituted.\par
The full higher-coherence regularization construction, all external sheaf-theoretic estimates, and the Bezrukavnikov–Yun calculations for every group have not been independently rebuilt. The inspected chain does not establish a failure of Theorems 3.4.8 or 4.4.1.\par
\section*{Outcome and remaining work}
The supported overall assessment is Major because the universal constructible pushforward conclusion has an admissible counterexample. P2–P5 have the limited impacts and corrections specified above. The main renormalized category and the Hecke equivalence are not disproved.\par
The publisher, arXiv, and author research records were searched for related source material, corrections, and supplementary code. Manuscript PDFs and TeX-source links were identified; no separate executable package or supplement for this exact file was identified. The author’s page warns that its hosted versions may differ from arXiv, so those files are not silently substituted. No author code was executed.\par
Item | Assessment | Verified scope\par
P1 | Major | Small-category pushforward fails for \(B\mathbb G_m\)\par
P2 | Minor | Adjunction strictness needs continuity\par
P3 | Minor | Galois sum correction preserves the needed retract\par
P4 | Minor | Descent requires a cover\par
P5 | Minor | Optional eigenvalue index needs inverse closure\par
Decision | Pending | Initial review; no independent or human acceptance\par
Delivery | Incomplete | Actual source/report image viewing missing\par
\section*{Supporting checks}
Paper ID: `M25-0360`.\par
> 00002TOKEN.\par
Per-problem evidence, raw source extractions, the reviewer script and actual output/exit receipt, coverage, provenance, dates, compiler/dependency payloads, and support inventory are retained under `verification/0360`.\par
This assistant session performs the scientific reasoning and report writing. Native tools perform mechanical algebra, compilation, parsing, and rendering; they are not independent reviewers. Infinite categorical proofs are not labelled as finite computational verification. Controller UI observations are attributed to the controller; backend equivalence is unverified. The image-reading interface returns `EMFILE`; rendering is not viewing, and completion remains withheld.\par
9 Q. P. Ho and P. Li, Revisiting mixed geometry, arXiv:2202.04833v4, submitted 1 December 2025; exact bound PDF internally dated 2 December 2025, https://arxiv.org/abs/2202.04833v4. EMS Press, DOI 10.4171/jems/1739, https://ems.press/journals/jems/articles/14299335, consulted 27 October 2026. Q. P. Ho, research publication list, https://quocho.com/research/, consulted 27 October 2026.\par
\end{document}
